\section{LASA 体系结构}
\label{sec:architecture}

\subsection{软硬件边界}

\cref{fig:system-placeholder} 给出完整推理系统。Cortex-A53 负责网络 DAG、特征
缓冲区和非卷积算子，四路 AXI DMA 分别传输输入、权重、偏置和输出；LASA
负责全部 13 个卷积。参数在一组图像处理期间常驻 DDR，层间张量保持紧凑 HWC
布局。该边界使通用处理器处理分支、池化和后处理等控制不规则工作，PL 则持续
执行规则且乘加密集的卷积。DMA 只观察连续字节区间，无需理解阵列行列或窗口
银行，从而把模型图调度与阵列内部映射分离。

\begin{figure*}[t]
\centering
\placeholderbox{52mm}{视觉占位：A53、DDR、四路 DMA、窗口生成、重放缓存、18$\times$16 双输出阵列、部分和反馈和 HWC 写回\\用粗箭头区分输入、参数、部分和及输出}
\caption{LASA 所在的边缘 MPSoC 推理系统。13 个卷积由 PL 执行，网络图和非卷积
算子由 A53 完成；外部张量始终保持紧凑 HWC 格式。}
\label{fig:system-placeholder}
\end{figure*}

每个卷积 dispatch 描述输入和输出形状、卷积核、空间 tile、输出通道块和量化
参数。LASA 按输入 tile、输出通道块和 $K$ 维轮次三个层次执行。输入 tile 只在
第一个层次进入窗口生成单元；输出通道块选择 32 个输出通道；$K$ 维轮次选择
当前 18 个规约元素。调度器改变三维遍历范围，数据通路规模保持不变。

\begin{table}[t]
\centering
\caption{LASA 片上存储的功能分工。}
\label{tab:memory-role}
\small
\begin{tabular}{lp{42mm}}
\toprule
存储 & 主要职责 \\
\midrule
URAM & 长层窗口向量、跨输出块重放的数据主体 \\
BRAM & 输入行、部分和、权重及局部数据队列 \\
LUTRAM & 窄控制队列、地址摘要和 HWC 字节打包 \\
寄存器 & 流水切片、有效位、上下文和提交状态 \\
\bottomrule
\end{tabular}
\end{table}

LASA 的数据通路可分为装载、变换、计算和提交四级。装载级以 AXI4-Stream
接收连续字节\cite{axi4stream}；变换级根据卷积核生成 18 元素向量；计算级
完成乘加和部分和反馈；提交级完成重量化、激活和 HWC 打包。每一级都具有独立
的容量和反压边界，因此输入突发、权重突发和输出写回可在不同周期推进。控制
信息只沿相邻级传播，避免形成贯穿整个加速器的组合就绪路径。

\subsection{从 HWC 字节到阵列向量}

3$\times$3 路径直接接收 HWC 字节。窗口生成单元使用行缓存保存相邻输入行，
按照输出坐标形成
\begin{equation}
v_{y,x,p}[r]=X[y+\Delta y,x+\Delta x,c],\quad 0\le r<18,
\label{eq:window-vector}
\end{equation}
其中 $r$ 与 $(\Delta y,\Delta x,c)$ 由当前轮次 $p$ 共同确定。图像边界产生量化
零值，最后一轮不足 18 个元素的位置由有效位屏蔽。生成的向量按输出像素和
轮次顺序写入片上缓存，同一地址序列随后可被不同输出通道块读取。

窗口缓存使用 URAM 保存长层的向量主体，BRAM 保存行缓存、地址信息和局部
队列，LUTRAM 用于窄而高频的控制和打包缓冲。存储分工由访问宽度和生命周期
决定：URAM 提供容量，BRAM 处理行级局部性，寄存器与 LUTRAM 隔离反压控制。
这种分工也避免宽数据寄存器使能直接由跨模块就绪链驱动，有利于维持
200 MHz 时序。

\begin{figure}[t]
\centering
\placeholderbox{44mm}{视觉占位：HWC 输入像素、三行缓存、18 元素分段和重放银行\\标出边界零值、尾段有效位和地址次序}
\caption{HWC 输入的窗口生成与银行映射。外部不产生完整 im2col，窗口向量在
片上按像素和计算轮次组织。}
\label{fig:bank-placeholder}
\end{figure}

1$\times$1 路径绕过三行窗口，只按通道方向每次提取 18 个输入。m14、m16 和
两个检测头使用该路径。它与 3$\times$3 路径共享缓存接口和阵列入口，因此
切换卷积核形态只改变向量生成方式。原生路径避免将 1$\times$1 权重映射到
稀疏 3$\times$3 窗口，也降低无效输入和轮次数量。

\subsection{18$\times$16 双输出通道阵列}

阵列由 18 行、16 列处理单元组成。每个处理单元在同一输入上并行计算两个输出
通道，因而总并行度为
\begin{equation}
N_{MAC}=18\times16\times2=576.
\label{eq:mac-count}
\end{equation}
输入沿规约方向分配，权重在当前上下文期间驻留，部分和通过 DSP48E2 级联在列
方向累加\cite{dsp48e2}。第 $p$ 轮输出为
\begin{equation}
s_p[c_o]=s_{p-1}[c_o]+\sum_{r=0}^{17}v_p[r]w_p[c_o,r],
\label{eq:pass-sum}
\end{equation}
其中首轮从偏置开始，末轮进入重量化，其他轮写入部分和缓冲。双输出结构把
32 个输出通道组织为同一块，既匹配 DMA 打包宽度，也使偏置和量化参数能够按
块装载。

每个双输出处理单元包含两个乘法累加通路，共享同一输入而读取两组权重。输入
广播只跨 16 列，部分和则沿 DSP 级联局部传递；这比把 32 个输出通道直接展开
为宽交叉连接更利于布局。18 行对应一个 $K$ 维轮次，尾轮通过行有效位保持
相同流水深度。阵列在一个上下文内只接受固定权重和有效位，权重切换发生在
上下文边界，因而数据周期不携带通道地址和模型信息。

当 $C_{out}$ 不是 32 的整数倍时，有效通道掩码同时作用于权重、偏置、重量化
和写回。两个 255 通道检测头的末块含 31 个有效通道。写回单元跨像素和 tile
维护字节余量，在最后一个 AXI beat 上产生精确 TKEEP，使 26$\times$26$\times$255
和 13$\times$13$\times$255 的紧凑张量不带填充字节。尾通道由此只损失一个
阵列 lane，不改变软件可见形状。

\subsection{部分和、激活与写回}

非末轮的 INT32 结果按像素和输出通道块写入部分和缓冲，并在下一轮按相同顺序
返回阵列。末轮结果加入偏置后按照层量化参数执行整数缩放、舍入和饱和，再经
256 项查找表实现激活\cite{jacob2018quant,leakyrelu2015}。写回单元将 32 lane
结果恢复为像素优先 HWC 字节流。由此，输入和输出均保持软件使用的布局，窗口
展开、规约分段和输出块只存在于 LASA 内部。

数据通路在各级使用 ready/valid 寄存切片。输入反压保持未消费字节，窗口缓存
只在完整向量提交后发布，部分和只在对应计算上下文确认后释放。错误状态会停止
新上下文并保留首个错误原因，使错误执行与正常完成具有不同的可见状态。该控制
原则为下一节的重叠调度提供确定边界。

这种边界还有一个物理实现含义：高扇出的上下文判断只驱动窄指针和状态，宽输入、
权重和部分和寄存器由局部已登记条件更新。最终 200 MHz 设计的关键路径因此
位于局部控制和存储使能附近，而不是从系统级描述直接扇出到数千位数据使能。
体系结构通过局部提交点把协议正确性与宽数据时序同时纳入设计。
